% Responsable: P4  (versión completa en extendido/05_set_de_prueba.tex)
\section{Conjunto de prueba y metodología}
\label{sec:set_de_prueba}

El conjunto (\texttt{data/test\_set.csv}) tiene 73 oraciones en tres niveles
(tabla~\ref{tab:set_de_prueba}). Cada una lleva la salida esperada: si la gramática debe
aceptarla, si es español correcto, el tipo de error, los fenómenos que prueba, la categoría de
cada palabra, el sujeto, el verbo, el objeto y a qué se une la primera frase preposicional.
\begin{itemize}
  \item \textbf{Nivel 1 (simples):} sujeto-verbo y sujeto-verbo-objeto, todas aceptables.
  \item \textbf{Nivel 2 (intermedias):} adjetivos, intensificadores, adverbios y frases
    preposicionales, más errores de concordancia (\emph{*La perro come}).
  \item \textbf{Nivel 3 (complejas):} varias frases preposicionales, ambigüedad estructural,
    concordancia a distancia y 11 oraciones correctas que la gramática no genera.
\end{itemize}

\begin{table}[h]
  \centering
  \small
  \resizebox{\linewidth}{!}{%
  \begin{tabular}{lrrrrrr}
    \toprule
    Nivel & Oraciones & En la gramática & Válidas fuera & Error sintáctico & Error de concordancia & Con PP \\
    \midrule
    1 & 20 & 20 & 0 & 0 & 0 & 0 \\
    2 & 29 & 18 & 0 & 1 & 10 & 6 \\
    3 & 24 & 7 & 11 & 14 & 3 & 11 \\
    Total & 73 & 45 & 11 & 15 & 13 & 17 \\
    \bottomrule
  \end{tabular}}
  \caption{Composición del conjunto de prueba.}
  \label{tab:set_de_prueba}
\end{table}

\paragraph{Dos referencias.} Cada decisión se evalúa contra \texttt{en\_gramatica} (¿sigue las
reglas?) y contra \texttt{espanol\_valido} (¿es español correcto?). Las 11 oraciones correctas
fuera de la gramática separan a un sistema que sigue las reglas de uno que sigue su conocimiento
del español.

\paragraph{Métricas.} Con la clase positiva \emph{acepta}:
precisión $= VP/(VP+FP)$, cobertura $= VP/(VP+FN)$, rechazo correcto $= VN/(VN+FP)$ y
exactitud. Además: tipo de error correcto en los rechazos, exactitud del etiquetado palabra por
palabra, exactitud de sujeto, verbo y objeto, y exactitud de la unión de la frase preposicional
(una falta de respuesta cuenta como error). Cada sistema participa solo donde produce datos.

\paragraph{Sistemas.} Parser propio; spaCy 3.8 (\texttt{es\_core\_news\_sm}); Transformers
(XLM-R para etiquetar y BETO para la gramaticalidad); y el LLM \texttt{gemini-3.5-flash-lite}
con temperatura~0, en versiones \emph{zero-shot} y \emph{few-shot}. Todos se corrieron el 26 de
noviembre de 2026; las métricas salen de \texttt{python -m evaluation.metrics}.
